% 加载数据
\DTLloaddb{mydata}{blind_review/publications.csv}


% 定义一个空白行的宏，方便后面循环调用 (保持原有格式)
\newcommand{\printemptyrow}{%
    \parbox[c][1.97cm][c]{0.432cm}{\raggedright\songti\zihao{5} }  &
    \parbox[c][1.97cm][c]{5.92cm}{\raggedright\songti\zihao{5} } &
    \parbox[c][1.97cm][c]{2.01cm}{\raggedright\songti\zihao{5} } &
    \parbox[c][1.97cm][c]{1.61cm}{\raggedright\songti\zihao{5} } &
    \parbox[c][1.97cm][c]{2.01cm}{\raggedright\songti\zihao{5} } &
    \parbox[c][1.97cm][c]{1.61cm}{\raggedright\songti\zihao{5} } \\
    \hline%
}
% 定义一个计数器来记录打印了多少行
\newcounter{rowprinted}

% --- 1. 定义单行渲染格式 (辅助宏) ---
% 这一步是为了防止在缓存扩展时出错，把格式封装起来
\newcommand{\renderMyRow}[6]{%
    \parbox[c][1.97cm][c]{0.432cm}{\centering\songti\zihao{5} #1} &
    \parbox[c][1.97cm][c]{5.92cm}{\raggedright\songti\zihao{5} #2} &
    \parbox[c][1.97cm][c]{2.01cm}{\raggedright\songti\zihao{5} #3} &
    \parbox[c][1.97cm][c]{1.61cm}{\centering\songti\zihao{5} #4} &
    \parbox[c][1.97cm][c]{2.01cm}{\centering\songti\zihao{5} #5} &
    \parbox[c][1.97cm][c]{1.61cm}{\centering\songti\zihao{5} #6} \\
    \hline
}
% --- 2. 准备缓存变量 ---
\newcommand{\tableContentBuffer}{} % 初始化为空

% --- 3. 在表格 *外部* 执行循环，构建缓存 ---
\setcounter{rowprinted}{0}

\DTLforeach{mydata}{%
    \id=序号, \pub=刊物名称, \auth=作者, \year=年份, \chap=章节, \idx=索引%
}{%
    % 使用 \xappto (Expand Append) 将当前行的数据填充进去
    % \noexpand 防止格式命令被提前展开，只展开 \id, \pub 等变量
    \xappto\tableContentBuffer{%
        \noexpand\renderMyRow{\id}{\pub}{\auth}{\year}{\chap}{\idx}%
    }%
    \stepcounter{rowprinted}%
}

% --- 4. 在表格 *外部* 补充空行 ---
\whiledo{\value{rowprinted} < 5}{%
    % 直接追加空行宏
    \gappto\tableContentBuffer{\printemptyrow}%
    \stepcounter{rowprinted}%
}

\cleardoublepage
{
  \heiti\zihao{-2}\bfseries{}
  \chapter{攻读硕士学位期间取得的研究成果}
}
\songti\zihao{5}
\linespread{2}\selectfont

{%
  \setlength{\parindent}{2.7em}
  \linespread{2.0}
  \indent 一、已发表（包括已录用待发表）的论文\par
}

\vspace{2mm}
{
  \centering\linespread{1.2}
  % \renewcommand{\arraystretch}{2.5}  % 控制行高，根据需要调整数值
  \begin{adjustbox}{width=\textwidth}
    \begin{tabular}{
      | >{\centering\arraybackslash}p{0.432cm}
      | >{\centering\arraybackslash}p{5.92cm}
      | >{\arraybackslash}p{2.01cm}
      | >{\centering\arraybackslash}p{1.61cm}
      | >{\arraybackslash}p{2.01cm}
      | >{\arraybackslash}p{1.61cm} |}
      \hline
      % 表头
      \parbox[c][2.16cm][c]{\linewidth}{\songti\zihao{5}\bfseries 序号} &
      \parbox[c][2.16cm][c]{\linewidth}{\songti\zihao{5}\bfseries 发表或投稿刊物/会议名称} &
      \parbox[c][2.16cm][c]{\linewidth}{\songti\zihao{5}\bfseries 作者（仅注明第几作者）} &
      \parbox[c][2.16cm][c]{\linewidth}{\songti\zihao{5}\bfseries 发表年份} &
      \parbox[c][2.16cm][c]{\linewidth}{\songti\zihao{5}\bfseries 与学位论文哪一部分（章、节）相关} &
      \parbox[c][2.16cm][c]{\linewidth}{\songti\zihao{5}\bfseries 被索引收录情况} \\
      \hline

      % 数据行
      \tableContentBuffer
      
    \end{tabular}
  \end{adjustbox}
}

\vspace{4mm}


{%
  \setlength{\parindent}{2.5em}
  \linespread{2.0}
  \indent 二、与学位论文内容相关的其它成果（包括专利、著作、获奖项目等）\par
  % \setlength{\hangindent}{2.5em} % 悬挂对齐到label
  % \setlength{\hangindent}{5.5em} % 悬挂对齐到label后
  \indent 专利：已授权一项发明专利，第一发明人，2020年1月\par

  % \setlength{\hangindent}{2.5em} % 悬挂对齐到label
  % \setlength{\hangindent}{5.5em} % 悬挂对齐到label后
  \indent 著作：参与编写一本著作，第二作者，2020年1月\par
  
  % \setlength{\hangindent}{2.5em} % 悬挂对齐到label
  % \setlength{\hangindent}{7.5em} % 悬挂对齐到label后
  \indent 获奖项目：获省部级科技奖一项，第三获奖人，2020年1月\par

  % \setlength{\hangindent}{2.5em} % 悬挂对齐到label
  % \setlength{\hangindent}{5.5em} % 悬挂对齐到label后
  \indent 竞赛：获数学建模竞赛全国一等奖，第一获奖人，2020年1月\par
}
